% !TeX root = manuscript.tex
% Related work draft, 2026-11-06. References use verified primary metadata.
\section{Related Work}
\label{sec:related_work}

\subsection{Visual Speech Generation and Automatic Dubbing}

Visual speech generation uses facial motion to constrain the content or timing of synthesized speech. Video-to-speech synthesis typically infers spoken content from silent video, whereas automatic dubbing additionally accepts a target script and may use reference speech to specify a voice. LipVoicer~\cite{yemini2024lipvoicer}, for example, combines a lip-reading model with diffusion-based speech generation and automatic-speech-recognition guidance. This combination illustrates how inferred linguistic content can complement visual conditioning. Script-conditioned dubbing places particular emphasis on satisfying the supplied text while matching the target articulation and reference voice.

Several dubbing systems address these requirements through explicit prosody and style modeling. V2C~\cite{chen2022v2c} studies visual voice cloning with text, visual information, and a reference voice. HPMDubbing~\cite{cong2023hpmdubbing} models prosody hierarchically using lip, face, and scene information, and associates linguistic and visual representations through text--lip attention followed by learned temporal upsampling. StyleDubber~\cite{cong2024styledubber} develops multiscale style modeling and uses forced-aligned phoneme information to support duration and reference-style processing. These methods highlight the value of visual and vocal cues at different granularities, although their alignment mechanisms should be distinguished: learned temporal upsampling and externally supervised phoneme timing impose different training requirements.

Neural codec language models provide another route to automatic dubbing. VoiceCraft-Dub~\cite{sungbin2025voicecraftdub} incorporates visual conditioning into a speech language model through visual adapters and feature fusion, and introduces CelebV-Dub for evaluating dubbing in expressive, real-world videos. Our work instead studies continuous acoustic generation with flow matching, focusing on the interaction of acoustic, visual, and textual representations. We also examine video-to-speech synthesis and visual forced alignment as related applications. The latter aims to localize a supplied transcript in video; it is distinct from using an external forced aligner to prepare dubbing training data.

\subsection{Implicit Alignment in Flow-Based Speech Generation}

Flow matching learns a time-dependent vector field that transports a simple distribution toward a target data distribution~\cite{lipman2023flowmatching}. In speech synthesis, F5-TTS~\cite{chen2025f5tts} demonstrates non-autoregressive generation without an explicit duration model or a separate phoneme-alignment stage. Extending this formulation to dubbing requires learning how linguistic units relate to visual articulation as well as to acoustic frames. CTC~\cite{graves2006ctc} provides sequence supervision without predefined token boundaries, but linguistic sequence consistency alone does not specify every aspect of audio--visual synchronization.

AlignDiT~\cite{choi2025aligndit} combines flow matching with CTC-based internal alignment for synchronized speech generation. It fuses acoustic inputs with resampled visual features by channel concatenation, and uses the resulting representation to query encoded text in its transformer blocks. Text thus serves as a conditioning sequence, while acoustic and visual information share the evolving hidden representation. AlignDiT also demonstrates extensions to video-to-speech synthesis and visual forced alignment. Our work builds on its implicit-alignment setting and investigates a different fusion mechanism: separate audio, video, and text states are updated through joint attention, with a soft temporal bias applied to audio--visual interactions.

CoSyncDiT~\cite{cong2026cosyncdit} explicitly controls the order in which conditions enter the generator. Its cognitive synchronous flow separates acoustic-style modeling, visual refinement through gated residual injection, and linguistic integration through time-adaptive cross-attention. Its Joint Semantic and Alignment Regularization (JSAR) combines CTC with frame-level contrastive representation alignment. This design addresses conditional interference through staged integration and auxiliary supervision. Our architecture instead refines the three modality streams jointly and encodes local temporal preference directly in the audio--visual attention logits.

UltraDub~\cite{cong2026ultradub} introduces visually steered flow learning and Rhythm-anchored Trajectory Guidance (RTG). Its text and speaker-style branches use queries that combine the evolving acoustic representation with lip features, while retrieving information from separate condition memories with independently normalized attention. During sampling, trajectory guidance evaluates progressively conditioned predictions at a visually guided midpoint to adjust the generation trajectory. Our approach focuses on interactions within the generator: audio, video, and text participate in jointly normalized attention, and visual and textual hidden states are updated alongside acoustic states. The Gaussian temporal bias specifies a local preference between audio and video frames rather than an inference-time trajectory correction.

\subsection{Multimodal Diffusion Transformers and Temporal Conditioning}

Diffusion transformers replace convolutional denoisers with transformer-based generative backbones~\cite{peebles2023dit}. MM-DiT extends this idea by maintaining modality-specific processing paths and allowing them to communicate through joint attention, as demonstrated in rectified-flow image generation~\cite{esser2024scaling}. This separation preserves distinct representations while permitting cross-modal exchange within each multimodal block. It provides the architectural basis for our generator; our contribution concerns its use for implicitly aligned dubbing and the temporal structure of its attention.

MMAudio~\cite{cheng2025mmaudio} applies joint audio, video, and text modeling to video-to-audio synthesis, followed by audio-only processing. It combines semantic visual features with higher-rate synchronization features and uses frame-aligned conditioning to guide sound generation. ImmersiveTTS~\cite{yun2026immersivetts} uses a multimodal diffusion transformer for environment-aware speech synthesis, together with global environmental conditioning and domain-specific representation alignment. Its content path uses alignment and duration modeling, and its environmental conditions serve a different role from the lip-motion sequence required for dubbing. These studies establish the usefulness of multimodal generative backbones while addressing different conditioning and temporal requirements.

For script-conditioned dubbing, audio and video share a sampled time axis, but text does not have known frame-level boundaries. Our design represents this distinction within joint attention: a Gaussian bias favors temporally neighboring audio--visual pairs in both directions, while within-modality and text-related attention remain global. This soft window does not impose a hard correspondence or remove attention connections outside a prescribed band. It supplies a local temporal prior that complements implicit linguistic alignment and joint multimodal refinement.
